\documentclass[11pt]{article}
\usepackage[margin=1in]{geometry}
\usepackage{amsmath,amssymb,amsthm}
\usepackage{booktabs}
\usepackage{graphicx}
\usepackage{array}
\usepackage{hyperref}
\usepackage{enumitem}
\usepackage{caption}
\usepackage{mathtools}
\captionsetup{font=small}

\newtheorem{theorem}{Theorem}
\newtheorem{proposition}[theorem]{Proposition}
\newtheorem{lemma}[theorem]{Lemma}
\newtheorem{definition}[theorem]{Definition}
\newtheorem{remark}[theorem]{Remark}

\DeclareMathOperator{\TIP}{TIP}
\DeclareMathOperator{\supp}{supp}
\DeclareMathOperator{\Rips}{Rips}
\DeclareMathOperator{\diam}{diam}
\DeclareMathOperator*{\argmin}{arg\,min}

\title{CC-SPR Supplementary Materials}
\author{}
\date{}

\begin{document}
\maketitle

\tableofcontents
\newpage

% ============================================================
\section{Mathematical Background}
\label{sec:supp_background}
% ============================================================

\subsection{Simplicial homology and persistence}

Let $X = \{x_1, \ldots, x_n\} \subset \mathbb{R}^p$ be a point cloud equipped with a metric $d$.
The \emph{Vietoris--Rips complex} at scale $\varepsilon$ is
\[
\Rips_\varepsilon(X, d) = \bigl\{ \sigma \subseteq X : \diam_d(\sigma) \le \varepsilon \bigr\}.
\]
As $\varepsilon$ increases from $0$ to $\infty$, the nested family $\{\Rips_\varepsilon\}_{\varepsilon \ge 0}$ forms a filtration.
Applying simplicial homology with field coefficients $\mathbb{F}$ to this filtration yields a persistence module that decomposes into interval modules $\{[b_i, d_i)\}_{i \in I}$.

\begin{definition}[Persistence barcode]\label{def:barcode}
The \emph{barcode} of a filtered simplicial complex is the multiset of intervals $\{[b_i, d_i)\}_{i \in I}$ encoding the birth and death scales of homology classes.
The \emph{persistence} of the $i$-th feature is $\pi_i = d_i - b_i$, and its \emph{midlife} is $m_i = (b_i + d_i)/2$.
\end{definition}

\begin{definition}[Chain groups and boundary maps]\label{def:chains}
Let $K$ be a simplicial complex.
The \emph{$k$-chain group} $C_k(K; \mathbb{F})$ is the $\mathbb{F}$-vector space with basis the $k$-simplices of $K$.
The \emph{boundary operator} $\partial_k: C_k \to C_{k-1}$ is defined on each $k$-simplex $\sigma = [v_0, \ldots, v_k]$ by
\[
\partial_k(\sigma) = \sum_{i=0}^k (-1)^i [v_0, \ldots, \hat{v}_i, \ldots, v_k].
\]
The fundamental identity $\partial_{k-1} \circ \partial_k = 0$ ensures $H_k(K; \mathbb{F}) = \ker(\partial_k) / \operatorname{im}(\partial_{k+1})$.
\end{definition}

\subsection{Harmonic representatives}

For a selected $H_1$ bar $[b, d)$, let $z_0$ be an initial cycle representative and $\partial_2$ the boundary operator from 2-chains to 1-chains.
The \emph{harmonic representative} minimizes $\ell_2$ energy within the affine homology class:
\[
\hat{z}_{\mathrm{harm}} = \argmin_{y \in C_2} \| z_0 + \partial_2 y \|_2^2.
\]

\begin{remark}[Dense support problem]\label{rem:dense}
The harmonic solution typically has $\|\hat{z}_{\mathrm{harm}}\|_0 = O(|E|)$ where $|E|$ is the number of edges.
In omics applications, this means hundreds or thousands of nonzero edge weights, diluting the biological signal when projected to feature space.
\end{remark}

% ============================================================
\section{Proofs and Derivations}
\label{sec:supp_proofs}
% ============================================================

\subsection{Convexity of the CC-SPR objective}

\begin{proposition}[Convexity]\label{prop:convex}
The CC-SPR objective
\[
f(y) = \| z_0 + \partial_2 y \|_2^2 + \lambda_1 \| z_0 + \partial_2 y \|_1 + \lambda_2 \| y \|_2^2
\]
is convex in $y$. When $\lambda_2 > 0$, it is strongly convex with a unique minimizer.
\end{proposition}

\begin{proof}
Each term is convex:
(i) $f_1(y) = \|z_0 + \partial_2 y\|_2^2$ is convex as a composition of a convex quadratic with an affine map;
(ii) $f_2(y) = \lambda_1 \|z_0 + \partial_2 y\|_1$ is convex (norms composed with affine maps);
(iii) $f_3(y) = \lambda_2 \|y\|_2^2$ is convex (strongly convex when $\lambda_2 > 0$).
A nonnegative sum of convex functions is convex, and strong convexity of any summand implies strong convexity of the sum.
\end{proof}

% ============================================================
\section{Backend Implementation Details}
\label{sec:supp_backends}
% ============================================================

\subsection{Euclidean backend}
Standard $\ell_2$ distance: $D_{ij} = \|x_i - x_j\|_2$.
Applied after optional PCA dimensionality reduction.

\subsection{Ollivier--Ricci curvature backend}
\begin{enumerate}
\item Construct $k$-nearest-neighbor graph $G = (V, E, w)$ with $w_{ij} = \|x_i - x_j\|_2$.
\item For $T$ iterations: compute Ollivier--Ricci curvature $\kappa(i,j)$ for each edge via optimal transport between neighbor distributions; update weights $w_{ij}^{(t+1)} = w_{ij}^{(t)} - \eta \cdot \kappa^{(t)}(i,j) \cdot w_{ij}^{(t)}$.
\item Compute shortest-path distances on the deformed graph.
\end{enumerate}
Default parameters: $k=10$, $T=5$, $\eta=0.5$.

\subsection{Diffusion backend}
\begin{enumerate}
\item Construct affinity matrix $W_{ij} = \exp(-\|x_i - x_j\|_2^2 / \sigma_i \sigma_j)$ where $\sigma_i$ is adaptive bandwidth.
\item Form transition matrix $P = D^{-1}W$ where $D = \text{diag}(W\mathbf{1})$.
\item Compute diffusion distance: $D_{ij}^{\text{diff}} = \|P^t_{i\cdot} - P^t_{j\cdot}\|_2$ at time scale $t$.
\end{enumerate}

\subsection{PHATE-like backend}
Compute potential distances from the diffusion operator:
$D_{ij}^{\text{pot}} = \|\log(P^t_{i\cdot} + \epsilon) - \log(P^t_{j\cdot} + \epsilon)\|_2$
with $\epsilon = 10^{-7}$ for numerical stability.

\subsection{Distance-to-measure (DTM) backend}
$d_i^{\text{DTM}} = \frac{1}{k}\sum_{j \in \text{kNN}(i)} \|x_i - x_j\|_2^2$, then
$D_{ij}^{\text{DTM}} = |d_i^{\text{DTM}} - d_j^{\text{DTM}}|$.
Robust to outliers: points in sparse regions have high DTM values, reducing their influence on the filtration.

% ============================================================
\section{Resource-Bounded Single-Cell Validation}
\label{sec:supp_scrna}
% ============================================================

Three GSE165087 configurations failed under cluster memory limits, establishing the resource boundary:

\begin{table}[h]
\centering
\small
\begin{tabular}{lllll}
\toprule
Run type & State & Elapsed & Requested memory & Peak RSS \\
\midrule
scRNA full & OUT\_OF\_MEMORY & 16:43:25 & 96\,GB & 100.3\,GB \\
scRNA fast-track & TIMEOUT & 20:00:06 & 96\,GB & 75.2\,GB \\
scRNA low-memory & OUT\_OF\_MEMORY & 11:25:59 & 48\,GB & 49.7\,GB \\
scRNA manuscript-safe & COMPLETED & 00:07:49 & 32\,GB & 2.9\,GB \\
scRNA ablation-safe & COMPLETED & --- & 32\,GB & ${\sim}$9.3\,GB \\
\bottomrule
\end{tabular}
\caption{Execution envelope for single-cell experiments on GSE165087.}
\label{tab:supp_scrna}
\end{table}

The manuscript-safe configuration ($n=320$ cells) produces F1 $= 0.925$ for all methods (ceiling effect).
The ablation-safe configuration ($n=400$) shows CC-SPR ($0.823$) exceeding harmonic ($0.798$) while trailing standard ($0.888$).

% ============================================================
\section{Ablation Tables}
\label{sec:supp_ablation}
% ============================================================

\begin{table}[h]
\centering
\small
\begin{tabular}{llccc}
\toprule
Dataset & Ablation axis & Best value & Mean F1 & $n$ \\
\midrule
LUAD & distance mode & ricci & 0.463 & 2 \\
LUAD & Ricci iterations & 10 & 0.551 & 2 \\
LUAD & $k$ & 10 & 0.405 & 2 \\
LUAD & $\lambda$ & $10^{-7}$ & 0.594 & 2 \\
LUAD & normalize & std & 0.522 & 2 \\
LUAD & top-$K$ & 10 & 0.667 & 2 \\
\midrule
GSE161711 & distance mode & euclidean & 0.600 & 5 \\
GSE161711 & Ricci iterations & 0 & 0.578 & 5 \\
GSE161711 & $k$ & 5 & 0.632 & 5 \\
GSE161711 & $\lambda$ & $10^{-4}$ & 0.590 & 5 \\
GSE161711 & normalize & std & 0.596 & 5 \\
GSE161711 & top-$K$ & 50 & 0.806 & 5 \\
\bottomrule
\end{tabular}
\caption{Best ablation settings from cached sweeps across LUAD and GSE161711.}
\label{tab:supp_ablation}
\end{table}

% ============================================================
\section{Full Enrichment Results}
\label{sec:supp_enrichment}
% ============================================================

Complete enrichment tables with all terms at adjusted $p < 0.1$ are provided in the supplementary CSV files.
The enrichment analysis was performed using gseapy v1.1.12 with ORA against MSigDB Hallmark 2020, KEGG 2021 Human, and Reactome 2022 databases using Enrichr's default background gene set.

% ============================================================
\section{Software Versions}
\label{sec:supp_software}
% ============================================================

\begin{table}[h]
\centering
\small
\begin{tabular}{ll}
\toprule
Package & Version \\
\midrule
Python & $\geq$ 3.10 \\
GUDHI & $\geq$ 3.8 \\
CVXPY & $\geq$ 1.4 \\
GraphRicciCurvature & $\geq$ 0.5 \\
scanpy & $\geq$ 1.9 \\
scikit-learn & $\geq$ 1.3 \\
gseapy & $\geq$ 1.1 \\
numpy & $\geq$ 1.24 \\
\bottomrule
\end{tabular}
\caption{Key software dependencies. A frozen environment file (\texttt{requirements-frozen.txt}) with exact versions is provided in the repository.}
\label{tab:supp_software}
\end{table}

\end{document}
